% ============================================================================
% 3. DW-BENCH DESIGN — Journal version (expanded with formal definitions)
% ============================================================================
\section{DW-Bench: Benchmark Design}
\label{sec:benchmark}

\subsection{Problem Formulation}

We formalize the data warehouse schema reasoning task as follows.
Given a heterogeneous schema graph $G = (V, E_{fk} \cup E_{lin})$
where $V$ is the set of tables, $E_{fk}$ is the set of foreign key
edges, and $E_{lin}$ is the set of data lineage (derived-from) edges,
and a natural language question $q$ about the topology of $G$, the
task is to produce an answer $a$ that matches the ground-truth answer
$a^* = f(G, q)$ computed by the corresponding graph algorithm $f$.

The ground-truth function $f$ is deterministic: for path queries,
$f$ returns the shortest path via BFS; for component queries, $f$
returns weakly connected components; for lineage queries, $f$ returns
directed reachability results.  This determinism ensures evaluation
integrity without human annotation.

\subsection{Datasets}

We select five datasets representing diverse data warehouse topologies:

% TABLE 1: Dataset Overview
\begin{table}[t]
\centering
\caption{DW-Bench dataset statistics. ``Silos'' = number of disconnected
components in the combined FK+lineage graph.}
\label{tab:datasets}
\begin{tabular}{lcccccc}
\toprule
Dataset & Domain & Tables & FK & Lineage & Questions & Silos \\
\midrule
AdventureWorks & Retail/HR    & 102 & 136 & 39 & 208 & 11 \\
TPC-DS         & Analytics    &  24 &  70 &  0 & 127 &  1 \\
TPC-DI         & ETL          &  35 &  29 & 21 & 181 &  2 \\
OMOP CDM       & Healthcare   &  37 &  74 & 21 & 158 &  3 \\
Syn-Logistics  & Supply Chain &  64 &  96 & 35 & 372 &  5 \\
\midrule
\textbf{Total} &              & \textbf{262} & \textbf{405} & \textbf{116} & \textbf{1,046} &: \\
\bottomrule
\end{tabular}
\end{table}

\textbf{AdventureWorks}~is a Microsoft reference data warehouse with 102
tables spanning OLTP and DW layers, connected by 136 FK edges and 39
lineage (derived\_from) edges.  Its dual-layer structure with 11 connected
components makes it ideal for lineage impact and silo detection questions.

\textbf{TPC-DS}~\cite{tpcds2021} is the industry-standard analytics benchmark with a star
schema of 24 tables and 70 FK edges.  Its single connected component and
absence of lineage edges tests pure FK-based reasoning in a well-studied
schema.

\textbf{TPC-DI}~is the data integration benchmark with 35 tables modeling
an ETL pipeline from staging to warehouse, with 21 lineage edges
representing the transformation flow from source to target tables.

\textbf{OMOP CDM}~(Observational Medical Outcomes Partnership Common Data
Model)~\cite{ohdsi2019omop} is a healthcare standard with 37 tables, 74
FK edges, and 21 lineage edges across 3 connected components (clinical
data, vocabulary tables, and metadata).

\textbf{Syn-Logistics}~is a synthetic supply-chain data warehouse with 64
tables modeled across 5 connected components (carrier, procurement,
healthcare, finance, HR).  It was specifically designed to address the
statistical power limitation of the real-world corpus: every subtype has
$n\geq20$ questions, enabling reliable per-subtype conclusions.

\subsection{Schema Graph Representation}

Each dataset is represented as a heterogeneous graph using PyTorch
Geometric~\cite{fey2019pyg} \texttt{HeteroData}.  Nodes represent
tables with 6 structural features: in-degree, out-degree, normalized
degree, lineage degree, betweenness centrality, and PageRank.  Edges are
typed as either \texttt{fk\_to} (foreign key) or \texttt{derived\_from}
(data lineage).  Lineage edges form a strict DAG (directed acyclic graph):
each \texttt{derived\_from} edge points from a downstream DW table to its
upstream source, with no cycles or self-references.  All traversal
algorithms respect edge typing: FK queries use only \texttt{fk\_to}
edges, lineage queries use only \texttt{derived\_from}, and
\texttt{combined\_impact} queries compose both types in two phases.

\paragraph{Definition 1 (\texttt{combined\_impact} semantics).}
\label{def:combined_impact}
Given a seed table $s$ and schema graph $G = (V, E_{\text{fk}}, E_{\text{lin}})$,
the \texttt{combined\_impact} answer is computed in two phases:
(1)~\emph{Lineage closure}: compute the set $L(s)$ of all tables reachable
from $s$ by following \texttt{derived\_from} edges forward (downstream
consumers of $s$);
(2)~\emph{FK dependents}: for each $t \in L(s) \cup \{s\}$, collect all tables
referencing $t$ via at least one \texttt{fk\_to} edge.
The answer is $L(s) \cup \bigcup_{t \in L(s) \cup \{s\}} \text{FK-deps}(t)$,
excluding $s$ itself.
This sequential composition (lineage then FK) reflects the enterprise
impact-analysis workflow: propagate the data change through the ETL pipeline,
then identify which downstream relational tables are invalidated.

\subsection{Question Generation}
\label{sec:qa-generation}

Questions are generated deterministically from the graph structure using
standard graph algorithms (BFS, Dijkstra, connected component detection)
implemented with NetworkX~\cite{networkx2008}, ensuring ground-truth
answers are provably correct.  We define three categories with 13 Tier-1
subtypes:

% TABLE 2: Question Taxonomy
\begin{table}[t]
\centering
\caption{Question taxonomy.  Difficulty is assigned based on the number of
reasoning hops and edge types required.}
\label{tab:taxonomy}
\small
\begin{tabular}{llclr}
\toprule
Category & Subtype & Difficulty & Description & \#Qs \\
\midrule
\multirow{5}{*}{\rotatebox{90}{\small Lineage}}
  & forward         & Easy & Direct lineage targets   & 72 \\
  & reverse         & Easy & Source tables for a DW table & 71 \\
  & transitive      & Hard & Multi-hop lineage chains  & 33 \\
  & combined\_impact & Hard & Lineage + FK dependents  & 69 \\
  & multi\_source   & Medium & Tables with 3+ sources   & 49 \\
\midrule
\multirow{3}{*}{\rotatebox{90}{\small Route}}
  & direct\_fk      & Easy & FK adjacency check       & 100 \\
  & join\_path      & Medium & FK shortest path         & 342 \\
  & hop\_count      & Medium & Path length              & 124 \\
\midrule
\multirow{5}{*}{\rotatebox{90}{\small Silo}}
  & count           & Easy & Number of components     & 30 \\
  & membership      & Hard & Which component contains X? & 38 \\
  & isolation       & Medium & Is table X isolated?     & 60 \\
  & connected       & Medium & Are X and Y connected?   & 29 \\
  & full\_enumeration & Hard & List all tables in a silo & 29 \\
\bottomrule
\end{tabular}
\end{table}

For Tier~2, we define 8 value-level subtypes that require interactive
access to row-level data: \texttt{cascade\_count},
\texttt{row\_provenance}, \texttt{row\_impact}, \texttt{value\_origin},
\texttt{multi\_hop\_trace}, \texttt{value\_propagation},
\texttt{cross\_silo\_reachability}, and \texttt{shared\_source}.

\subsection{Difficulty Assignment}

Questions are assigned three difficulty levels:
\begin{itemize}
  \item \textbf{Easy}: single-hop queries (direct FK adjacency, forward/reverse
    lineage) requiring only one edge traversal.
  \item \textbf{Medium}: multi-hop queries within a single edge type
    (FK shortest path, hop count, isolation check, component connectivity).
  \item \textbf{Hard}: queries requiring composition of multiple edge types
    or exhaustive enumeration (combined impact, transitive lineage,
    membership, full enumeration).
\end{itemize}

\subsection{Obfuscation Protocol}

To distinguish genuine topology reasoning from schema name memorization,
we generate an \emph{obfuscated} variant for each dataset.  Table names
are replaced with random identifiers (\texttt{Table\_A},
\texttt{Table\_B}, etc.) using a deterministic mapping.  Questions and
answers are updated accordingly with word-boundary-aware replacement to
avoid corrupting natural language phrases.
